\documentclass[11pt]{article}
\usepackage{fullpage,graphicx,hyperref,amsmath,amssymb}

\newcommand{\fig}[2]{\begin{figure}[h]\begin{center}%
  \includegraphics[scale=#2]{#1.pdf}\end{center}%
  \end{figure}}
\newcommand{\ceil}[1]{\left\lceil #1\right\rceil}

\begin{document}

\begin{center}\Large\bf 
CS 473 (Spring 2025)\\
Homework 10 (due May 1 Thu 10am)
\end{center}

\medskip\noindent
{\bf Instructions:} As in previous homeworks.

\begin{description}



\bigskip
\item[Problem 10.1:]
Consider the following problem:
The input consists of $n$ rectangles, where the $i$-th rectangle $R_i$
has height $h_i\le 1$, width $w_i\le 1$, and area $a_i=h_iw_i$.
The goal is to place the rectangles in a rectangular container of width 1
while minimizing the height of the 
container.  Rotations of the rectangles are not allowed.

Consider the following greedy algorithm: 
\begin{quote}
\begin{tabbing}
1.M\= for\ \=\kill
0.\> sort the rectangles so that $h_1\ge h_2\ge\cdots\ge h_n$\\
1.\> $k_1=1$\\
2.\> for $m=1,2,\ldots$ until $k_m=n+1$ \{\\
3.\>\> let $k_{m+1}$ be the largest index such that
$w_{k_m}+w_{k_m+1}+\cdots+w_{k_{m+1}-1}\le 1$\\
4.\>\> place $R_{k_m},R_{k_m+1},\ldots, R_{k_{m+1}-1}$ in a new row (a ``shelf'') of
height $h_{k_m}$\\
\} 
\end{tabbing}
\end{quote}

Show that this greedy algorithm runs in polynomial time and
has approximation factor at most 3.

[Hint: Prove the inequality
$a_{k_m}+a_{k_m+1}+\cdots+a_{k_{m+1}} > h_{k_{m+1}}$.
Then sum both sides\ldots]

\newpage
\textbf{Solution:}
\paragraph{Polynomial running time.}
Sorting the $n$ rectangles by non‐increasing height takes $O(n\log n)$ time.  
Then the single pass that scans through the list, forming each shelf by accumulating widths until they exceed 1, clearly takes $O(n)$ time.  Hence the overall running time is $O(n\log n)$.

\paragraph{Approximation ratio.}
Let $H_G$ be the total height of the packing produced by the greedy algorithm, and let $H^*$ be the height of an optimal packing.

Label the shelves produced by the algorithm $m=1,2,\dots,M$, and for each $m$ let $k_{m+1}$ be the index of the first rectangle not placed on shelf $m$, so by construction 
\[
w_{k_m}+w_{k_m+1}+\cdots+w_{k_{m+1}-1} \le 1
\quad\text{but}\quad
w_{k_m}+w_{k_m+1}+\cdots+w_{k_{m+1}} > 1.
\]
Each shelf $m$ is given height $h_{k_m}$, so
\[
H_G  = \sum_{m=1}^M h_{k_m}\,.
\]
For each shelf $m$,
\[
\sum_{i=k_m}^{k_{m+1}} a_i
= \sum_{i=k_m}^{k_{m+1}} h_i\,w_i
\ge h_{k_{m+1}} \bigl(w_{k_m}+ \cdots + w_{k_{m+1}}\bigr)
> h_{k_{m+1}}.
\]
Summing this strict inequality over $m=1,2,\dots,M$ gives
\[
\sum_{m=1}^M \sum_{i=k_m}^{k_{m+1}} a_i
> \sum_{m=1}^M h_{k_{m+1}}
= \sum_{m=2}^{M+1}h_{k_m}
= \sum_{m=2}^M h_{k_m}.
\]
(Since $k_{M+1}=n+1$ and we may set $h_{n+1}=0$.) The left-hand side is
\[
\sum_{m=1}^M \sum_{i=k_m}^{k_{m+1}} a_i 
= \sum_{m=1}^M \left(\sum_{i=k_m}^{k_{m+1}-1} a_i + a_{k_{m+1}}\right)
= A + \sum_{m=1}^M a_{k_{m+1}}
< 2A
\]
and in any packing into a width-1 strip of height $H^*$ the total area satisfies $A\le H^*$.  Hence
\[
\sum_{m=2}^M h_{k_m}
< 2A \le 2H^*.
\]
Therefore
\[
H_G
= \sum_{m=1}^M h_{k_m}
= h_{k_1} + \sum_{m=2}^M h_{k_m}
< h_{k_1}  +  2H^*.
\]
Finally, since the tallest rectangle has height $h_{k_1}$ and must fit into any feasible packing, we have
\[
h_{k_1} \le H^*,
\]
and thus
\[
H_G  
< h_{k_1}+H^*
\le 3 H^*.
\]
The approximation factor is at most 3.

\newpage
\item[Problem 10.2:]
Consider the following variant of the traveling salesman problem: given a complete
weighted graph $G=(V,E)$, where the weights satisfy $d(u,v)=d(v,u)$
and $d(u,v)\le d(u,w)+d(w,v)$, find a cycle $C$ that
visits every vertex exactly once, minimizing the 
{\em maximum edge weight\/} in $C$.  (In other words, our salesman
doesn't care about the total travel time, but can't tolerate
long flights and wants to keep each flight short.) 
Equivalently, we seek the smallest weight $r^*$
such that $G[r^*]$ contains a Hamiltonian cycle,
where $G[r]$ denotes the subgraph $(V,\{uv\in E\mid d(u,v)\le r\})$.

Consider the following approach:
\begin{quote}
\begin{tabbing}
1.M \=\kill
1.\> find the smallest $r_0$ such that $G[r_0]$ is connected\\
2.\> let $T$ be any spanning tree of $G[r_0]$\\
3.\> get a (simple) cycle $C$ from $T$ and return $C$
\end{tabbing}
\end{quote}
To perform line~3, we pick an arbitrary node~$u$
as the root of $T$ and call the following recursive procedure,
which returns a simple path in $G$ that visits all descendants of $u$:
\begin{quote}
\begin{tabbing}
1.M \=\kill
visit($u$):\\
1.\> if $u$ is a leaf then return $\langle u\rangle$\\
2.\> let $u_1,\ldots,u_k$ be the children of $u$\\
3.\> for $i=1,\ldots,k$ do $\pi_i=\mbox{visit}(u_i)$\\
4.\> return $\langle u\rangle \circ (\pi_1)^R \circ (\pi_2)^R \circ\cdots \circ (\pi_k)^R$
\end{tabbing}
\end{quote}
Here, $\circ$ denotes concatenation, and $\pi^R$ denotes the reverse
of the list~$\pi$.  We get~$C$ by completing
the resulting path into a cycle in $G$.  

\begin{enumerate}
\item[(a)] (70 pts) Show that the cycle $C$ produced by the above procedure has
maximum edge weight $\le 3r_0$.  

[Hint: You may want to first deduce some
properties about the path generated by visit().  Try a few examples
and see; then use induction to prove your observations formally.]

\medskip
\item[(b)] (30 pts) Conclude that there is a polynomial-time algorithm with
approximation factor at most 3.  
\end{enumerate}

\newpage
\textbf{Solution:}

\begin{enumerate}
\item[(a)]
\textit{Lemma.} For any node $u$, the path $P(u) = \langle p_0, p_1,\cdots,p_m \rangle$ produced by \texttt{visit}$(u)$ satisfies:
\begin{enumerate}
    \item The last node $p_m$ is one of the children of $u$.
    \item The first edge $(p_0, p_1)$ satisfies $d(u, p_1) \le 2r_0$.
    \item The closing edge $(p_m, u)$ satisfies $d(p_m, u) \le r_0$.
\end{enumerate}

\textit{Proof.} We prove by induction on the structure of $T$.

\textbf{Base case:} If $u$ is a leaf, $P(u) = \langle u \rangle$, and all three properties hold trivially.

\textbf{Inductive step:} Suppose for each child $u_i$, the path $P(u_i)$ satisfies the lemma. Then:
\begin{itemize}
  \item The last node $p_m$ is the first node of $(P(u_k))^R$, which is $u_k$, a child of $u$.
  \item The first node after $u$ is the last node of $(P(u_1))^R$, which is the last node of $P(u_1)$, say $q_{m_1}$. By inductive hypothesis, $q_{m_1}$ is a descendant of $u_1$. By the triangle inequality:
  \[
  d(u, q_{m_1}) \le d(u, u_1) + d(u_1, q_{m_1}) \le r_0 + r_0 = 2r_0.
  \]
  \item The closing edge from $p_m$ back to $u$ connects $u_k$ to $u$, and since $u_k$ is a child of $u$, the edge has weight at most $r_0$.
\end{itemize}
This completes the inductive proof of the lemma. $\hfill\square$

\medskip

Now consider the edges in the final cycle $C$:
\begin{itemize}
  \item \textbf{Internal edges within each $(P(u_i))^R$:}  
  These edges are either edges within a child subtree or its reverse. Since each subtree path corresponds to traversals in $T$, and every edge in $T$ has weight at most $r_0$, these internal edges are of weight at most $r_0$.
  
  \item \textbf{Edges between subpaths:}  
  When transitioning from the end of $(P(u_i))^R$ to the beginning of $(P(u_{i+1}))^R$, the two nodes involved are $u_i$ (the child just completed) and some descendant $q_{m_{i+1}}$ of $u_{i+1}$. By the triangle inequality:
  \[
  d(u_i, q_{m_{i+1}}) \le d(u_i, u) + d(u, q_{m_{i+1}}).
  \]
  Since $u_i$ and $u$ are connected by an edge in $T$ of weight at most $r_0$, and from the lemma $d(u, q_{m_{i+1}}) \le 2r_0$, we get:
  \[
  d(u_i, q_{m_{i+1}}) \le r_0 + 2r_0 = 3r_0.
  \]
  
  \item \textbf{Closing edge:}  
  Finally, the edge from $p_m$ (which is $u_k$) back to $u$ has weight at most $r_0$ by the lemma.
\end{itemize}

Thus, every edge in the constructed cycle $C$ has weight at most $3r_0$.

\item[(b)]
The algorithm provided in the question yields a polynomial-time, 3-approximation algorithm.

\begin{enumerate}
  \item[1.] \textbf{Computing \(r_0\).}  
  We sort all \(\binom{n}{2}\) edge weights in non-decreasing order and scan them, maintaining the connectivity of the subgraph \(G[r]\) using a Union-Find data structure. This allows us to determine the smallest \(r_0\) such that \(G[r_0]\) is connected. The time complexity of this step is \(O(n^2 \log n)\).

  \item[2.] \textbf{Building and traversing the spanning tree.}  
  Once \(r_0\) is known, we construct any spanning tree \(T\) of \(G[r_0]\) using Prim's algorithm in \(O(n^2)\) time. The recursive \texttt{visit} procedure then generates a path through all nodes in linear time relative to the number of vertices.

  \item[3.] \textbf{Approximation.}  
  As shown in (a), the final Hamiltonian cycle \(C\) has maximum edge weight at most \(3r_0\). Since any valid Hamiltonian cycle in \(G\) must use only edges of weight \(\le r^*\), where \(r^*\) is the optimal value, and \(r_0 \le r^*\), we conclude:
  \[
      \max_{e \in C} d(e) \le 3r_0 \le 3r^*.
  \]
\end{enumerate}

Therefore, the algorithm runs in polynomial time and achieves a 3-approximation factor for the problem.


\end{enumerate}



\newpage
\item[Problem 10.3:]
Given a {\em directed\/} graph $G=(V,E)$, we want to find
a subgraph $A$ that is acyclic (i.e., a dag), maximizing the
number of edges in $A$. 
\begin{enumerate}
\item[(a)] (50 pts) Show that the following deterministic algorithm yields
a feasible solution with approximation factor at least $1/2$:
\begin{quote}
\begin{tabbing}
1.M\= for\ \=\kill
1.\> fix an arbitrary (not random) order of the vertices $v_1,\ldots, v_n$\\
2.\> for $i=2$ to $n$ do \{\\
3.\>\> $\textrm{IN}_i =$ all edges from $\{v_1,\ldots,v_{i-1}\}$ to $v_i$\\
4.\>\> $\textrm{OUT}_i =$ all edges from $v_i$ to $\{v_1,\ldots,v_{i-1}\}$\\
5.\>\> if $|\textrm{IN}_i|\ge |\textrm{OUT}_i|$ then insert $\textrm{IN}_i$ to $A$ else
insert $\textrm{OUT}_i$ to $A$\\
\>\}
\end{tabbing}
\end{quote}

\medskip
\item[(b)] (50 pts) Show that the following randomized algorithm yields
a feasible solution and analyze its expected approximation factor:
\begin{quote}
\begin{tabbing}
1.M\= for\=\kill
1.\> take a random order of the vertices $v_1,\ldots,v_n$\\
2.\> for each edge $(v_i,v_j)\in E$ do\\
3.\>\> if $i<j$ then insert $(v_i,v_j)$ to $A$
\end{tabbing}
\end{quote}
\end{enumerate}

[Note: it may be useful to recall the fact that a directed graph is a dag iff there exists a topological sort.]

\newpage
\textbf{Solution:}
\newcommand{\IN}{\textrm{IN}}
\newcommand{\OUT}{\textrm{OUT}}
\newcommand{\E}{\mathbb{E}}
\begin{enumerate}
\item[(a)]
\textbf{Feasibility.}
By construction, every edge inserted into \(A\) goes between \(v_i\) and some
vertex in \(\{v_1,\dots,v_{i-1}\}\).  If we orient all chosen edges so that they go
forward in the order \(v_1\prec v_2\prec\cdots\prec v_n\), then no edge ever
points from a later vertex to an earlier one.  Hence \(A\) admits the topological
order \(v_1,\dots,v_n\) and is acyclic.

\textbf{\(1/2\)-Approximation.}
Let \(A\) be the edge set returned by the algorithm, and let \(E^*\) be an optimal
maximum-cardinality acyclic subgraph of \(G\).

Observe that every edge \((u,v)\) with \(u\neq v\) appears exactly once among
the union of all \(\IN_i\cup\OUT_i\), namely when the later of \(u,v\)
is being processed.  Hence
\[
\sum_{i=2}^n \bigl(|\IN_i| + |\OUT_i|\bigr)  =  |E|.
\]
By the choice in the algorithm, the edges added at step \(i\) is
\[
\max\{\,|\IN_i|, |\OUT_i|\,\}
\ge \frac12\bigl(|\IN_i|+|\OUT_i|\bigr).
\]
Summing over \(i\) gives
\[
|A|
= \sum_{i=2}^n \max\{|\IN_i|,|\OUT_i|\}
\ge \frac12 \sum_{i=2}^n \bigl(|\IN_i|+|\OUT_i|\bigr)
= \frac12\,|E|.
\]
Therefore
\[
|A| \ge \frac12\,|E|
\ge \frac12\,|E^*|,
\]
as claimed.

\item[(b)]

\textbf{Feasibility.}
Let \(\pi\) be the random order of vertices \((v_1, \dots, v_n)\).
By construction, every edge in \(A\) goes forward in the random order \(\pi\).  Hence \(\pi\) itself is a valid topological sort of the subgraph \((V,A)\), so \((V,A)\) is a DAG.

\textbf{Expected approximation ratio.}
Let \(E^*\subseteq E\) be an optimal maximum-cardinality acyclic subgraph, so \(|E^*| = \max\{|A'| : (V,A')\text{ is a DAG}\}\).  Note that trivially \(|E^*|\le |E|\).

For any fixed edge \(e=(u\to v)\in E\), under a uniformly random order \(\pi\) we have
\[
\Pr\bigl[e\text{ is kept in }A\bigr]
=\Pr\bigl(u\text{ precedes }v \text{ in }\pi\bigr)
=\frac12.
\]
Define the indicator random variable
\[
X_e =
\begin{cases}
1,&\text{if }e\in A,\\
0,&\text{otherwise.}
\end{cases}
\]
Then \(\E[X_e]=\frac12\), and linearity of expectation gives
\[
\E\bigl[\,|A|\,\bigr]
=\E\!\Bigl[\sum_{e\in E}X_e\Bigr]
=\sum_{e\in E}\E[X_e]
=\sum_{e\in E}\frac12
=\frac12\,|E|.
\]
Since \(|E|\ge|E^*|\), it follows that
\[
\E\bigl[\,|A|\,\bigr]
=\frac12\,|E|
\ge \frac12\,|E^*|.
\]
The expected approximation factor is \(\frac12\).

\end{enumerate}

\end{description}
\end{document}


